% Generated draft: requires mathematical and source review.
\documentclass[11pt]{article}
\usepackage[margin=25mm]{geometry}
\usepackage{fontspec}
\setmainfont{Latin Modern Roman}
\usepackage{amsmath,amssymb,mathtools,unicode-math}
\setmathfont{Latin Modern Math}
\usepackage{enumitem,needspace,xurl,hyperref}
\setlength{\parindent}{0pt}
\setlength{\parskip}{5pt}
\setlength{\emergencystretch}{3em}
\newcommand{\N}{\mathbb{N}}
\newcommand{\Z}{\mathbb{Z}}
\newcommand{\Q}{\mathbb{Q}}
\newcommand{\R}{\mathbb{R}}
\newcommand{\C}{\mathbb{C}}
\newcommand{\F}{\mathbb{F}}
\newcommand{\A}{\mathbb{A}}
\newcommand{\PP}{\mathbb{P}}
\newcommand{\E}{\mathbb{E}}
\newcommand{\eps}{\varepsilon}
\newcommand{\dd}{\,\mathrm{d}}
\begin{document}
\section*{UnsolvedMath ID: 30005740}
\textbf{Source problem code:} catalogue entry
\subsection{Definitions and mathematical statement}
Let \(\F\) be a field of characteristic \(0\), and let \(h(x)\in\F[x]\). Define the first Weyl algebra and the subalgebra appearing in the supplied extract by
\[
A_1=\F\langle x,y\rangle/(yx-xy-1),\qquad
A_h=\F\langle x,h(x)y\rangle\subseteq A_1.
\]
Here the second expression denotes the unital subalgebra generated by the indicated elements. The extract allows an arbitrary polynomial \(h\); no restriction excluding \(h=0\) is supplied.

For an associative \(\F\)-algebra \(A\), put \(A^e=A\otimes_{\F}A^{\mathrm{op}}\) and
\[
HH^n(A)=\operatorname{Ext}_{A^e}^n(A,A),\qquad
HH^\bullet(A)=\bigoplus_{n\geq0}HH^n(A).
\]
Use the canonical cup product \(\smile\) and Gerstenhaber bracket \([-,-]_G\), of degrees \(0\) and \(-1\), respectively. A Batalin--Vilkovisky structure compatible with these operations means an \(\F\)-linear operator
\[
\Delta:HH^n(A)\longrightarrow HH^{n-1}(A),\qquad \Delta^2=0,
\]
where \(HH^{-1}(A)=0\), such that, for all homogeneous \(a\in HH^p(A)\) and \(b\in HH^q(A)\),
\[
[a,b]_G=(-1)^p\bigl(\Delta(a\smile b)-\Delta(a)\smile b-(-1)^p a\smile\Delta(b)\bigr).
\]
This fixes a sign convention for the bracket-generating operator.

The supplied reconstructed statement additionally assumes that \(A_h\) is twisted Calabi--Yau. In the usual terminology, this means that \(A_h\) has a bounded resolution by finitely generated projective \(A_h^e\)-modules and, for some integer \(d\geq0\) and automorphism \(\nu\),
\[
\operatorname{Ext}_{A_h^e}^{i}(A_h,A_h^e)\cong
\begin{cases}
0,&i\ne d,\\
{}_1(A_h)_{\nu},&i=d,
\end{cases}
\]
as bimodules, using the convention that the right action on \({}_1(A_h)_{\nu}\) is \(u\cdot a=u\nu(a)\). Such an automorphism is called a Nakayama automorphism.

\textbf{Question, subject to the qualifications below.} For \(A_h\) as above, under the reconstructed twisted Calabi--Yau hypothesis and the source condition that its Nakayama automorphism is not semisimple, does \(HH^\bullet(A_h)\) admit an operator \(\Delta\) satisfying the displayed conditions?

\textbf{Unresolved qualifications.} The extract specifies the Weyl subalgebra family, whereas the reconstructed statement describes \(A_h\) only as a twisted Calabi--Yau algebra. It is therefore unclear whether the intended scope is precisely this family or all twisted Calabi--Yau algebras. The extract does not supply the additional twisted Calabi--Yau hypothesis or its applicable range of \(h\). It also does not specify the meaning of semisimplicity for the automorphism on this potentially infinite-dimensional algebra, the field over which it is to be tested, or whether the question seeks existence for every qualifying \(h\), for some qualifying \(h\), or a classification. The displayed question records the conditional existence problem for a qualifying member without resolving those scope issues. Compatibility with the canonical Gerstenhaber structure is the interpretation used here; the extract does not explicitly define its intended BV structure.
\subsection{Short English statement}
For the algebra generated by \(x\) and \(h(x)y\) inside the first Weyl algebra over a characteristic-zero field, ask whether its Hochschild cohomology has a square-zero operator of degree \(-1\) generating the usual Gerstenhaber bracket when the Nakayama automorphism is nonsemisimple. The supplied reconstruction adds a twisted Calabi--Yau assumption. Its scope, the precise semisimplicity condition, and the quantification over \(h\) require clarification.
\subsection{Sources}
\begin{itemize}
\item Local UnsolvedMath registry, permanent ID \(30005740\), source problem code \texttt{catalogue entry}, title \emph{BV Structures with Nonsemisimple Nakayama Automorphisms}. The supplied record belongs to the Oberwolfach Reports Open Problems collection; no individual proposer is named.
\item Supplied report citation: \emph{Mini-Workshop: Poisson and Poisson-type algebras} (2024), Oberwolfach Reports, pp.~2681--2715, DOI \texttt{10.4171/owr/2023/46}, \url{https://doi.org/10.4171/owr/2023/46}. The supplied literature-source label is \emph{Oberwolfach Report 46/2023}. The extract includes the page marker 2703. The record gives proposed year 2023; the difference between that year and the citation's 2024 label has not been resolved here.
\item Literature assessment dated \texttt{2026-08-22}: the supplied curation labels the problem open and states that BV structures are established for twisted Calabi--Yau or Frobenius algebras with semisimple Nakayama automorphism, while the nonsemisimple \(A_h\) case in the extract remains unresolved. This is supplied dated curation, not an independently verified literature assessment. Its listed source is \url{https://doi.org/10.4171/owr/2023/46}.
\item The record labels its statement \texttt{corrected\_verified} and says it was reconstructed and checked against the official report. These are supplied provenance claims; no independent source visit or verification is asserted here.
\item The extract attributes the cup product and Gerstenhaber bracket to Gerstenhaber via its reference [2], but supplies no bibliographic details for [2]. No additional citation is inferred.
\end{itemize}
\end{document}
